% !TEX root = ../main.tex
\chapter{Các công trình liên quan}
\label{Chapter2}

Chương này trình bày cơ sở lý thuyết và các công trình liên quan theo bốn nhóm.
Mục~\ref{sec:cf-truyen-thong} giới thiệu các phương pháp lọc cộng tác và những hạn chế
chính của chúng. Mục~\ref{sec:llm-rec} trình bày hai hướng khai thác LLM trong hệ thống
tư vấn và vấn đề độ dài ngữ cảnh. Mục~\ref{sec:cach-noi-cf-llm} tổng hợp các cơ chế tích
hợp biểu diễn cộng tác vào LLM, từ ánh xạ MLP, học tương phản và mã hóa văn bản đến thích
nghi tham số và Q-Former. Cuối cùng, Mục~\ref{sec:nen-tang-qformer} trình bày cơ sở kiến
trúc Q-Former, làm nền tảng cho phương pháp được đề xuất ở Chương~\ref{Chapter3}.

\section{Lọc cộng tác truyền thống}
\label{sec:cf-truyen-thong}

Lọc cộng tác (collaborative filtering -- CF) khai thác các mẫu tương tác giữa người dùng và
sản phẩm để ước lượng sở thích. Trực giác cơ bản là những người dùng có hành vi tương tự
trong quá khứ có xu hướng quan tâm đến các sản phẩm tương tự trong tương lai. Nhóm phương
pháp này chủ yếu dựa trên ma trận tương tác người dùng--sản phẩm và không yêu cầu thông
tin nội dung của sản phẩm~\cite{zhang2019deep}.

\subsection{Phân rã ma trận}
\label{ssec:mf}

Phân rã ma trận (matrix factorization -- MF)~\cite{mf2009} là một phương pháp nền tảng của
lọc cộng tác. MF ánh xạ mỗi người dùng $u$ và sản phẩm $i$ thành các véc-tơ ẩn $e_u, e_i \in
\mathbb{R}^d$, sao cho tích vô hướng $e_u^\top e_i$ phản ánh mức độ tương thích giữa người
dùng và sản phẩm. Tùy dạng dữ liệu, mô hình có thể được tối ưu bằng hồi quy trên các đánh
giá quan sát được hoặc bằng các mục tiêu xếp hạng cặp như BPR~\cite{bpr2009}.

Ưu điểm của MF là cấu trúc đơn giản, khả năng mở rộng tốt và hiệu quả đã được kiểm chứng
trên nhiều bộ dữ liệu. Tuy nhiên, phép tương tác dựa trên tích vô hướng giới hạn khả năng mô
hình hóa các quan hệ phi tuyến phức tạp. Neural Collaborative Filtering (NCF)~\cite{ncf2017}
mở rộng ý tưởng này bằng cách sử dụng mạng nơ-ron nhiều lớp để mô hình hóa tương tác
người dùng--sản phẩm linh hoạt hơn.

\subsection{Lọc cộng tác dựa trên đồ thị}
\label{ssec:lightgcn}

Dữ liệu tương tác người dùng--sản phẩm có thể được biểu diễn dưới dạng đồ thị hai phía.
LightGCN~\cite{lightgcn2020} khai thác cấu trúc này bằng cách lan truyền và tổng hợp tuyến
tính embedding giữa các nút lân cận qua nhiều lớp. Nhờ tận dụng quan hệ bậc cao trong đồ
thị tương tác, LightGCN đạt hiệu quả cao trên nhiều bộ dữ liệu gợi ý và trở thành một baseline
phổ biến cho các phương pháp dựa trên tín hiệu cộng tác.

\subsection{Hạn chế của lọc cộng tác thuần túy}
\label{ssec:han-che-cf}

Mặc dù có hiệu quả tốt trong các kịch bản giàu tương tác, lọc cộng tác truyền thống vẫn có ba
hạn chế đáng chú ý. \textbf{Thứ nhất, dữ liệu thưa:} phần lớn cặp người dùng--sản phẩm không
có tương tác quan sát được, khiến embedding của các thực thể ít tương tác kém tin cậy.
\textbf{Thứ hai, cold-start:} người dùng hoặc sản phẩm mới gần như không có tín hiệu cộng tác
để xây dựng biểu diễn. \textbf{Thứ ba, thiếu thông tin ngữ nghĩa:} embedding cộng tác chủ yếu
phản ánh cấu trúc đồng xuất hiện trong dữ liệu tương tác và không trực tiếp biểu diễn nội dung
văn bản của sản phẩm hay mô tả sở thích người dùng.

Các hạn chế này là một trong những động lực thúc đẩy việc kết hợp mô hình ngôn ngữ với hệ
thống tư vấn.

\section{Mô hình ngôn ngữ lớn cho hệ thống tư vấn}
\label{sec:llm-rec}

Nhờ được tiền huấn luyện trên tập dữ liệu văn bản quy mô lớn, LLM có khả năng biểu diễn
ngữ nghĩa và suy luận vượt trội so với các mô hình chỉ khai thác tín hiệu tương tác. Do đó,
việc tích hợp LLM vào hệ thống tư vấn đã trở thành một hướng nghiên cứu đáng chú ý trong
những năm gần đây~\cite{llmrec2023}.

\subsection{Hai hướng khai thác LLM}
\label{ssec:hai-huong}

Các phương pháp kết hợp LLM với hệ thống tư vấn có thể được chia thành hai nhóm chính.

\textbf{LLM-enhanced.} LLM được sử dụng như một bộ mã hóa để tạo biểu diễn ngữ nghĩa cho
người dùng, sản phẩm hoặc nội dung liên quan. Các biểu diễn này sau đó được đưa vào mô
hình gợi ý hoặc mô hình xếp hạng truyền thống. Trong thiết kế này, LLM chủ yếu đóng vai trò
trích xuất đặc trưng, còn quyết định gợi ý cuối cùng vẫn do mô hình phía sau thực hiện.

\textbf{LLMs-as-RS.} Bài toán gợi ý được biểu diễn trực tiếp dưới dạng tác vụ ngôn ngữ.
P5~\cite{p52022} là một công trình tiêu biểu, thống nhất nhiều nhiệm vụ gợi ý thành các cặp
câu hỏi--trả lời và huấn luyện một mô hình ngôn ngữ để xử lý chúng. BIGRec~\cite{bigrec2023}
tinh chỉnh LLM để sinh ra mô tả của một sản phẩm phù hợp với sở thích người dùng, sau đó
đối sánh biểu diễn của đầu ra sinh với biểu diễn của các sản phẩm thực tế trong hệ thống để
xác định sản phẩm được gợi ý.

Hướng LLMs-as-RS có lợi thế trong việc tận dụng tri thức ngữ nghĩa và hỗ trợ các trường hợp
cold-start thông qua mô tả văn bản. Luận văn này tiếp tục theo hướng đó nhưng tập trung vào
cách đưa tín hiệu cộng tác vào LLM một cách hiệu quả.

\subsection{Giới hạn độ dài ngữ cảnh}
\label{ssec:rao-can-context}

Một cách tự nhiên để cung cấp thông tin hành vi cho LLM là chuyển lịch sử tương tác thành
văn bản và đưa trực tiếp vào prompt. Tuy nhiên, lịch sử của một người dùng thực tế có thể gồm
hàng trăm hoặc hàng nghìn tương tác. Khi mỗi sản phẩm chiếm nhiều token, toàn bộ lịch sử
có thể nhanh chóng vượt quá giới hạn ngữ cảnh hoặc làm chi phí suy luận tăng đáng kể.

CoLLM~\cite{collm2023} chỉ ra rằng việc biểu diễn toàn bộ lịch sử dưới dạng văn bản không
phải lúc nào cũng khả thi. Một giải pháp là sử dụng biểu diễn cộng tác nén, được học từ lịch
sử tương tác, rồi chuyển biểu diễn này thành soft token để đưa vào LLM. Cách tiếp cận đó dẫn
tới bài toán trung tâm của luận văn: căn chỉnh biểu diễn cộng tác với không gian embedding
của LLM như thế nào để thông tin được khai thác hiệu quả.

\section{Tích hợp biểu diễn cộng tác vào LLM}
\label{sec:cach-noi-cf-llm}

Các công trình gần đây đề xuất nhiều cơ chế khác nhau để kết hợp tín hiệu cộng tác với LLM.
Phần này tập trung vào năm hướng có liên quan trực tiếp đến luận văn và phân tích vai trò, ưu
điểm cũng như hạn chế của từng cơ chế.

\subsection{Ánh xạ MLP: CoLLM}
\label{ssec:collm}

CoLLM~\cite{collm2023} là baseline quan trọng của luận văn. Mô hình sử dụng một MLP hai
lớp theo cấu trúc Linear--ReLU--Linear, với tầng ẩn có kích thước gấp mười lần chiều
embedding cộng tác, để ánh xạ biểu diễn cộng tác sang không gian embedding của LLM. Mỗi
véc-tơ cộng tác được chuyển thành một soft token và chèn vào prompt. LLM sau đó thực hiện
dự đoán nhị phân dưới dạng xác suất cho câu trả lời ``Yes'' hoặc ``No''.

Quy trình huấn luyện của CoLLM gồm hai bước chính. Ở bước đầu, LoRA được tinh chỉnh trên
phần prompt chỉ chứa thông tin văn bản để LLM thích nghi với nhiệm vụ gợi ý mà chưa sử
dụng tín hiệu cộng tác. Ở bước sau, mô-đun CIE được huấn luyện để ánh xạ biểu diễn cộng tác
của người dùng và sản phẩm vào không gian embedding của LLM, trong khi các thành phần đã
học ở bước trước được giữ cố định theo cấu hình mặc định.

Ưu điểm của CoLLM nằm ở cấu trúc đơn giản và chi phí bổ sung tương đối nhỏ. Tuy nhiên,
phép ánh xạ MLP được áp dụng độc lập trên từng embedding. Cùng một hàm ánh xạ được sử
dụng cho mọi mẫu và không trực tiếp tổng hợp quan hệ giữa nhiều nguồn tín hiệu cộng tác
trước khi chúng đi vào LLM. Đồng thời, mỗi véc-tơ cộng tác được cô đọng thành một soft token
duy nhất, nên mô-đun không có cơ chế riêng để phân bổ hoặc chọn lọc các thành phần thông
tin theo mức độ liên quan với ngữ cảnh hiện tại.

\subsection{Căn chỉnh tương phản: SellaRec}
\label{ssec:sellarec}

SellaRec~\cite{sellarec2024} căn chỉnh không gian cộng tác và không gian ngôn ngữ bằng học
tương phản (contrastive learning), sử dụng mục tiêu InfoNCE để kéo gần biểu diễn tương ứng
của cùng một sản phẩm và đẩy xa các biểu diễn không tương ứng. Cơ chế này cung cấp một
mục tiêu huấn luyện trực tiếp cho việc thu hẹp khoảng cách giữa hai không gian biểu diễn.

SellaRec đã được đánh giá trên bài toán CTR với AUC và uAUC trên MovieLens-1M, do đó
cung cấp bằng chứng thực nghiệm cho hướng căn chỉnh dựa trên học tương phản. Trong khi đó,
Q-Former đại diện cho một cơ chế khác: thay vì căn chỉnh chủ yếu thông qua một mục tiêu
tương phản ở mức biểu diễn, Q-Former sử dụng các truy vấn học được để chọn lọc và tổng hợp
tín hiệu cộng tác trước khi chuyển sang không gian của LLM. Vì vậy, SellaRec và CoQLLM có
thể được xem là hai hướng tiếp cận bổ sung về cơ chế hơn là các biến thể trực tiếp của nhau.

Cần lưu ý rằng SellaRec sử dụng Qwen2-7B-Base~\cite{qwen2}, trong khi luận văn sử dụng
Vicuna-7B-v1.5~\cite{vicuna2023}. Sự khác biệt về backbone và thiết lập dữ liệu hạn chế khả
năng diễn giải chênh lệch số liệu như một so sánh thuần túy giữa hai cơ chế căn chỉnh.

\subsection{Mã hóa văn bản: BinLLM}
\label{ssec:binllm}

BinLLM~\cite{binllm2024} chuyển biểu diễn cộng tác thành chuỗi ký tự nhị phân, sau đó đưa
chuỗi này vào LLM như văn bản thông thường. Thiết kế này có ưu điểm là đơn giản về giao
diện với LLM và không yêu cầu cơ chế soft token chuyên biệt.

Hạn chế của phương pháp nằm ở bước rời rạc hóa. Việc chuyển một véc-tơ liên tục thành chuỗi
nhị phân có thể làm mất một phần thông tin của biểu diễn ban đầu. Bên cạnh đó, chuỗi nhị
phân không có ngữ nghĩa tự nhiên trong từ vựng của LLM, khiến mô hình phải học thêm cách
diễn giải cấu trúc mã hóa này trước khi có thể khai thác nó cho nhiệm vụ gợi ý.

\subsection{Thích nghi tham số LLM: CoRA}
\label{ssec:cora-rel}

CoRA~\cite{cora2024} đưa tín hiệu cộng tác vào LLM theo một hướng khác. Thay vì chèn tín
hiệu vào không gian đầu vào dưới dạng soft token, CoRA biến đổi biểu diễn cộng tác thành các
cập nhật trọng số hạng thấp phụ thuộc từng mẫu và bổ sung chúng vào các lớp attention của
LLM. Nhờ đó, tính toán bên trong LLM có thể được điều chỉnh theo từng cặp người dùng--sản
phẩm.

Cơ chế này đại diện cho một hướng thiết kế khác với CoQLLM. CoRA tác động trực tiếp lên
tham số hiệu dụng của LLM trong quá trình suy luận, trong khi CoQLLM giữ đường tích hợp ở
không gian đầu vào thông qua soft token. Luận văn không hiện thực lại cơ chế sinh trọng số của
CoRA mà sử dụng kết quả của mô hình như một baseline bổ sung trên Amazon-Book để định
vị CoQLLM so với một hướng tích hợp khác.

\subsection{Căn chỉnh qua Q-Former: ILM}
\label{ssec:ilm}

ILM~\cite{ilm2024} là một công trình sử dụng Q-Former để kết nối thông tin cộng tác với mô
hình ngôn ngữ trong bài toán gợi ý. Q-Former sử dụng tập truy vấn học được kết hợp với
cross-attention để khai thác thông tin từ các biểu diễn cộng tác, thay vì áp dụng một phép chiếu
độc lập trên từng véc-tơ.

Cơ chế này phù hợp với mục tiêu căn chỉnh của luận văn vì cho phép mô hình học cách tổng
hợp tín hiệu cộng tác trước khi đưa vào LLM. Tuy nhiên, ILM không báo cáo kết quả theo thiết
lập CTR với AUC và uAUC như CoLLM, nên chưa cung cấp cơ sở để so sánh trực tiếp hiệu quả
của Q-Former với MLP trong cùng giao thức đánh giá.

\subsection{So sánh các phương pháp}
\label{ssec:tong-hop-so-sanh}

Bảng~\ref{tab:so-sanh-phuong-phap} tổng hợp các cơ chế tích hợp tín hiệu cộng tác được thảo
luận ở trên.

\begin{table}[ht]
\caption{So sánh các phương pháp tích hợp biểu diễn cộng tác vào LLM}
\label{tab:so-sanh-phuong-phap}
\begin{center}
\begin{tabular}{|L{2cm}|L{3cm}|L{3.5cm}|L{3.5cm}|}
\hline
\textbf{Phương pháp} & \textbf{Cơ chế tích hợp} & \textbf{Điểm mạnh} & \textbf{Hạn chế chính} \\
\hline
CoLLM~\cite{collm2023} & MLP pointwise, mỗi embedding thành một soft token & Đơn giản; baseline mạnh & Ánh xạ độc lập; thiếu cơ chế tổng hợp và chọn lọc thích nghi \\
\hline
SellaRec~\cite{sellarec2024} & Học tương phản & Căn chỉnh trực tiếp hai không gian biểu diễn & Khác backbone; khó quy chênh lệch số liệu hoàn toàn cho cơ chế \\
\hline
BinLLM~\cite{binllm2024} & Mã hóa nhị phân & Giao diện đơn giản với LLM; không cần soft token & Có thể mất thông tin khi rời rạc hóa \\
\hline
CoRA~\cite{cora2024} & Sinh cập nhật trọng số theo mẫu & Cá nhân hóa sâu trong tính toán của LLM & Can thiệp vào trọng số thay vì không gian đầu vào \\
\hline
ILM~\cite{ilm2024} & Q-Former và cross-attention & Có khả năng chọn lọc và tổng hợp tín hiệu & Chưa báo cáo AUC/uAUC theo thiết lập CoLLM \\
\hline
\end{tabular}
\end{center}
\end{table}

Tổng hợp các công trình cho thấy Q-Former là một hướng đáng khảo sát cho bài toán căn chỉnh
biểu diễn cộng tác với LLM nhờ khả năng tổng hợp tín hiệu thông qua các truy vấn học được.
Tuy nhiên, hướng này vẫn thiếu đánh giá trực tiếp trong thiết lập CTR với AUC/uAUC và chưa
được khảo sát đầy đủ cùng các cơ chế huấn luyện hướng tới đặc tính của hai độ đo này. Đây là
khoảng trống mà luận văn tập trung giải quyết.

\section{Cơ sở kiến trúc Q-Former}
\label{sec:nen-tang-qformer}

BLIP-2~\cite{blip2_2023} giới thiệu Q-Former (Querying Transformer) như một mô-đun trung
gian giữa bộ mã hóa ảnh đóng băng và mô hình ngôn ngữ đóng băng. Mục tiêu của Q-Former
là trích xuất một tập biểu diễn nhỏ gọn từ đặc trưng đầu vào liên tục, sau đó chuyển các biểu
diễn này sang dạng mà LLM có thể khai thác mà không cần cập nhật toàn bộ hai mô hình lớn.

Q-Former sử dụng một tập truy vấn học được $\mathbf{Q} \in \mathbb{R}^{N_q \times d}$, thường
gồm một số lượng hữu hạn query token, kết hợp với các lớp Transformer. Trong quá trình xử lý,
hai cơ chế attention đóng vai trò chính:

\begin{enumerate}
    \item \textit{Self-attention} giữa các truy vấn, cho phép các query token trao đổi thông tin và
    phối hợp trong việc phân chia vai trò biểu diễn.
    \item \textit{Cross-attention} từ truy vấn đến đặc trưng đầu vào, cho phép mỗi truy vấn tổng
    hợp thông tin từ nguồn đầu vào với trọng số phụ thuộc vào nội dung của mẫu.
\end{enumerate}

Đầu ra của Q-Former là một tập véc-tơ truy vấn đã được cập nhật, ký hiệu $\mathbf{Z} \in
\mathbb{R}^{N_q \times d}$. Trong BLIP-2, các véc-tơ này được chiếu sang không gian
embedding của LLM và ghép với prompt văn bản.

BLIP-2 huấn luyện Q-Former theo hai giai đoạn. Giai đoạn đầu tập trung học biểu diễn qua các
mục tiêu tương phản ảnh--văn bản (ITC), so khớp ảnh--văn bản (ITM) và sinh văn bản có điều
kiện (ITG). Giai đoạn sau căn chỉnh đầu ra của Q-Former với mô hình ngôn ngữ để các biểu
diễn trung gian có thể được LLM khai thác trong quá trình sinh.

Đối với bài toán của luận văn, đặc trưng ảnh được thay thế bằng biểu diễn cộng tác. Khi đó,
Q-Former được sử dụng như một mô-đun căn chỉnh: các truy vấn học được tổng hợp thông tin
từ embedding người dùng, sản phẩm mục tiêu và lịch sử trước khi đầu ra được chuyển sang
không gian embedding của LLM.

\section{Khoảng trống nghiên cứu và định hướng đề xuất}
\label{sec:tong-ket-ch2}

Từ các nhóm công trình đã phân tích, có thể rút ra hai nhận xét. Thứ nhất, việc đưa tín hiệu
cộng tác vào LLM dưới dạng biểu diễn nén là cần thiết khi lịch sử tương tác quá dài để biểu
diễn đầy đủ bằng văn bản. Thứ hai, các cơ chế tích hợp hiện có khác nhau đáng kể về cách căn
chỉnh giữa không gian cộng tác và không gian ngôn ngữ.

CoLLM cung cấp một baseline đơn giản dựa trên MLP pointwise, trong khi Q-Former cung cấp
khả năng tổng hợp nhiều nguồn tín hiệu thông qua truy vấn học được và attention.
ILM cho thấy hướng Q-Former có thể áp dụng trong gợi ý, nhưng chưa đánh giá theo AUC/uAUC
trong thiết lập CTR tương tự CoLLM. Do đó, luận văn tập trung vào việc xây dựng và đánh giá
một mô-đun căn chỉnh Q-Former trong cùng nhóm bài toán, đồng thời khảo sát các cơ chế huấn
luyện hướng tới AUC và uAUC.
